\documentclass[a4paper,11pt]{article}

\newcommand{\TPnumero}{Lab 3 -- Solution}
% =========================================================
% Tutorial / lab preamble — Time Series Analysis module
% article class + fancyhdr + tcolorbox + listings (English)
% Author : Aymen Ben Brik (aymen.benbrik@esprit.tn)
% =========================================================

\usepackage[utf8]{inputenc}
\usepackage[T1]{fontenc}
\usepackage[english]{babel}
\usepackage{lmodern}
\usepackage{amsmath, amssymb, amsthm, mathtools, bm}
\usepackage{graphicx}
\usepackage{booktabs}
\usepackage{array}
\usepackage{multirow}
\usepackage{colortbl}
\usepackage{xcolor}
\usepackage{tikz}
\usetikzlibrary{positioning, arrows.meta, calc, shapes, fit, backgrounds, matrix}
\usepackage{listings}
\usepackage[most]{tcolorbox}
\usepackage{geometry}
\geometry{a4paper, margin=2cm, headheight=15pt}
\usepackage{fancyhdr}
\usepackage[hidelinks]{hyperref}
\usepackage{enumitem}

% --- Colours ---
\definecolor{espritBlue}{RGB}{30, 60, 110}
\definecolor{espritAccent}{RGB}{200, 80, 30}
\definecolor{boxEx}{RGB}{30, 60, 110}
\definecolor{boxQ}{RGB}{180, 120, 0}
\definecolor{boxC}{RGB}{30, 100, 60}
\definecolor{boxAtt}{RGB}{200, 80, 30}

% --- Listings ---
\definecolor{codebg}{RGB}{248,248,248}
\definecolor{codeframe}{RGB}{220,220,220}
\definecolor{keyword}{RGB}{0,82,155}
\definecolor{comment}{RGB}{88,110,117}
\definecolor{string}{RGB}{133,153,0}

\lstdefinestyle{pythonstyle}{
  language=Python,
  basicstyle=\ttfamily\small,
  keywordstyle=\color{keyword}\bfseries,
  commentstyle=\color{comment}\itshape,
  stringstyle=\color{string},
  backgroundcolor=\color{codebg},
  frame=single,
  rulecolor=\color{codeframe},
  frameround=tttt,
  breaklines=true,
  showstringspaces=false,
  tabsize=2
}
\lstdefinestyle{Rstyle}{
  language=R,
  basicstyle=\ttfamily\small,
  keywordstyle=\color{keyword}\bfseries,
  commentstyle=\color{comment}\itshape,
  stringstyle=\color{string},
  backgroundcolor=\color{codebg},
  frame=single,
  rulecolor=\color{codeframe},
  frameround=tttt,
  breaklines=true,
  showstringspaces=false,
  tabsize=2
}
\lstset{style=pythonstyle}

% --- Common macros (configurable path) ---
\providecommand{\macrosPath}{../../preamble/macros.tex}
\input{\macrosPath}

% --- Exercise counter ---
\newcounter{excounter}
\renewcommand{\theexcounter}{\arabic{excounter}}

\newtcolorbox[use counter=excounter]{exercise}[1][]{
  enhanced, breakable,
  colback=boxEx!5, colframe=boxEx, fonttitle=\bfseries,
  title={Exercise~\theexcounter\ifstrempty{#1}{}{ -- #1}},
  arc=2pt, boxrule=0.6pt
}
\newtcolorbox{question}[1][]{
  enhanced, breakable,
  colback=boxQ!5, colframe=boxQ, fonttitle=\bfseries,
  title={Question\ifstrempty{#1}{}{ -- #1}},
  arc=2pt, boxrule=0.5pt
}
\newtcolorbox{solution}[1][]{
  enhanced, breakable,
  colback=boxC!5, colframe=boxC, fonttitle=\bfseries,
  title={Solution\ifstrempty{#1}{}{ -- #1}},
  arc=2pt, boxrule=0.5pt
}
\newtcolorbox{warning}[1][]{
  enhanced, breakable,
  colback=boxAtt!5, colframe=boxAtt, fonttitle=\bfseries,
  title={Warning\ifstrempty{#1}{}{ -- #1}},
  arc=2pt, boxrule=0.5pt
}

% --- Headers / footers ---
\pagestyle{fancy}
\fancyhf{}
\fancyhead[L]{\textsc{\moduleNom} -- \auteurNom}
\fancyhead[R]{\TPnumero}
\fancyfoot[C]{\thepage}
\fancyfoot[L]{\auteurInstitution}
\fancyfoot[R]{\auteurEmail}
\renewcommand{\headrulewidth}{0.4pt}
\renewcommand{\footrulewidth}{0.2pt}

% --- Placeholder ---
\providecommand{\TPnumero}{Lab}


\title{Lab 3: ACF / PACF\\\large \emph{Reference solution}}
\author{\auteurNom\\\auteurEmail\\\auteurInstitution}
\date{\anneeUniversitaire}

\begin{document}
\maketitle

\noindent\emph{Reference Python code in \texttt{correction.py}.}

\setcounter{excounter}{0}
\begin{exercise}[White noise]\end{exercise}

\begin{solution}
On 500 i.i.d.\ $\mathcal{N}(0, 1)$ draws (seed 42), the empirical ACF
spikes at $\abs{\widehat\rho(h)}$ are essentially all within
$1.96 / \sqrt{500} \approx 0.088$. Out of 30 lags tested,
typically $1$ or $2$ exceed the band — close to the expected
$0.05 \times 30 = 1.5$ false positives.
\end{solution}

\begin{exercise}[AR(1) simulation]\end{exercise}

\begin{solution}
Empirical ACF agrees well with $\rho(h) = \phi^h$ for $\phi = 0.3, 0.7,
-0.7$. For $\phi = 0.95$, the empirical ACF \emph{decays slower} than
the theoretical exponential due to finite-sample bias: with $T = 500$
the series is close to a random walk, and $\widehat\rho$ often
overestimates persistence at high lags.

The empirical PACF for $\phi = 0.7$ shows $\widehat\alpha(1) \approx 0.7$
and $\widehat\alpha(h) \approx 0$ for $h \geq 2$ — the cutoff signature
of an AR(1).
\end{solution}

\begin{exercise}[MA(1) simulation]\end{exercise}

\begin{solution}
Theoretical $\rho(1) = \theta / (1 + \theta^2)$:
\begin{itemize}
  \item $\theta = 0.5$: $\rho(1) = 0.5 / 1.25 = 0.4$.
  \item $\theta = -0.5$: $\rho(1) = -0.4$.
  \item $\theta = 0.9$: $\rho(1) = 0.9 / 1.81 \approx 0.497$.
\end{itemize}
Empirical ACF on $T = 500$ samples confirms these values to within
$\pm 0.05$. PACF decays geometrically, never cuts off — opposite
signature to AR(1).
\end{solution}

\begin{exercise}[Atlanta ACF/PACF]\end{exercise}

\begin{solution}
\textbf{(1)} The raw ACF of the Atlanta series shows strong, slowly
decaying autocorrelations at lags $12, 24, 36, \dots$ (annual
seasonality) and high values at small lags from the remaining trend.
The series is far from white noise.

\textbf{(2)} After decomposition (centred MA $k = 6$ + seasonal
averaging) the residual standard deviation is $\approx 2.88^\circ$F.

\textbf{(3)} The residual ACF is much cleaner: the seasonal spikes
disappear, but a few short-lag autocorrelations remain (the
residuals are approximately AR(1) with $\phi \approx 0.3$). The
Atlanta residual is \emph{not exactly white noise} — there is short-term
weather persistence to capture with an ARMA model (Chapter 5).
\end{solution}

\begin{exercise}[Ljung--Box]\end{exercise}

\begin{solution}
On the Atlanta residual, the Ljung--Box statistic at $m = 12$ is
$\approx 200$ with $p$-value $\approx 0$: \textbf{reject} the
white-noise hypothesis. The same conclusion holds at $m = 24, 36$.

On the simulated white noise (Exercise 1), $Q_{\text{LB}}(m=12)
\approx 6$--$15$ with $p$-value $> 0.10$: \textbf{do not reject} —
correct behaviour under $H_0$.
\end{solution}

\end{document}
